\documentclass[11pt,a4paper]{article}

\usepackage[margin=1in]{geometry}
\usepackage{setspace}
\doublespacing

\usepackage{amsmath,amssymb,amsfonts}
\usepackage{graphicx}
\usepackage{booktabs,longtable,array,colortbl,tabularx}
\usepackage{caption,float}
\usepackage[numbers,sort&compress]{natbib}
\usepackage[colorlinks=true,linkcolor=blue,citecolor=blue,urlcolor=blue]{hyperref}
\usepackage{lineno}
\linenumbers

% Fix line numbers for AMS math environments
\newcommand*\patchAmsMathEnvironmentForLineno[1]{%
  \expandafter\let\csname old#1\expandafter\endcsname\csname #1\endcsname
  \expandafter\let\csname oldend#1\expandafter\endcsname\csname end#1\endcsname
  \renewenvironment{#1}%
     {\linenomath\csname old#1\endcsname}%
     {\csname oldend#1\endcsname\endlinenomath}}%
\newcommand*\patchBothAmsMathEnvironmentsForLineno[1]{%
  \patchAmsMathEnvironmentForLineno{#1}%
  \patchAmsMathEnvironmentForLineno{#1*}}%
\AtBeginDocument{%
\patchBothAmsMathEnvironmentsForLineno{equation}%
\patchBothAmsMathEnvironmentsForLineno{align}%
\patchBothAmsMathEnvironmentsForLineno{flalign}%
\patchBothAmsMathEnvironmentsForLineno{alignat}%
\patchBothAmsMathEnvironmentsForLineno{gather}%
\patchBothAmsMathEnvironmentsForLineno{multline}%
}

\title{\textbf{Physician Exodus from Tunisia: Dynamic Workforce Projections, Retention Policy Counterfactuals, and Economic Evaluation to 2050}\\\large A Multi-State Continuous-Time Markov Simulation and Cost-Consequence Analysis}

\author{
  \textbf{Health Workforce and Health Economics Research Consortium for Tunisia}\\[0.5em]
  \small Department of Health Economics and Public Health Policy, Faculty of Medicine of Tunis, University of Tunis El Manar, Tunis, Tunisia\\
  \small Health System Sustainability and Human Resources for Health Collaborative, Tunis, Tunisia
}

\date{August 20, 2026}

\begin{document}

\maketitle

\begin{abstract}
\noindent \textbf{Background:} The escalating emigration of physicians from Tunisia poses a structural threat to the sustainability of the public health system, universal health coverage, and regional equity. Since 2021, annual medical departures have exceeded 1,450 doctors per year ($\sim$10\% of the active workforce), with over 6,000 physicians emigrating between 2021 and 2024.\\
\textbf{Methods:} We developed an empirically grounded, 36-cell discrete-time Markov cohort simulation model with a continuous-time competing-risk hazard matrix engine, integrating the WHO Health Labour Market Framework, discrete-choice econometrics, and CHEERS 2022 standards. The baseline active stock at $t_0 = 2020$ was initialized via 2D Iterative Proportional Fitting (IPF) on official census data, adjusting for inactive registrations ($\alpha = 0.88$). Survey intentions were scaled by specialty-specific realization factors ($\kappa \in [0.25, 0.49]$) and corrected for odds ratio non-collapsibility. Two scalar calibration parameters were estimated on a 2020--2021 training sample, validated on a 2022--2024 holdout sample, and cross-checked against French (DREES/RPPS) and German (Bundes{\"a}rztekammer) destination registries ($r = 0.871$). We simulated six policy counterfactuals (2025--2050): Status Quo, Public Remuneration Reform (+30\%), Hospital Infrastructure Modernization, Regional Rural Retention Package, Training Quota Expansion (+15\%), and Comprehensive Synergistic Reform.\\
\textbf{Results:} Under Status Quo, active domestic physician density collapses from 11.4 per 10,000 population in 2020 to 7.7 per 10k by 2050, while the cumulative emigrated diaspora reaches 23,767 physicians abroad ($\sim$2.2$\times$ the domestic active workforce). Expanding training quotas alone (SCEN\_4) exacerbates brain drain (+2,303 additional emigrants), creating 117 Million TND in unrecovered public educational deadweight loss. Hospital infrastructure modernization (SCEN\_2) is the most cost-effective single lever ($\text{ICER}_{\text{public}} = 209,137\text{ TND}$ per discounted public physician-year retained). The Comprehensive Reform package (SCEN\_5) maximizes retention, generating 10,040 discounted public physician-years (15,151 undiscounted) and 4,058 interior physician-years at an $\text{ICER}_{\text{public}}$ of 210,374 TND (discounted net cost of 2,112.1 Million TND, averaging $\sim$2.8\% of the annual health budget). In 1,000-draw Monte Carlo PSA, SCEN\_5 achieves $>$80\% cost-effectiveness probability at willingness-to-pay thresholds above 250,000 TND/physician-year.\\
\textbf{Conclusion:} Supply expansion without retention fails to resolve health workforce depletion. A synergistic, multi-attribute policy package combining public salary adjustment, hospital modernization, and regional hardship incentives is macro-fiscally feasible and essential to safeguard Tunisia's health system through 2050.\\
\textbf{Keywords:} Physician Migration; Health Workforce Planning; Tunisia; Markov Simulation Model; Cost-Consequence Analysis; Discrete Choice Experiment; Human Resources for Health; Universal Health Coverage.
\end{abstract}

\newpage

\section{Introduction} \label{sec:introduction}

\subsection{Regional Context: Physician Emigration Across North African Health Systems}

The escalating emigration of highly trained medical professionals from low- and middle-income countries (LMICs) represents a profound structural challenge to the sustainability of national health systems, the achievement of Universal Health Coverage (UHC), and global health equity \citep{who2021, sousa2023, oecd2019}. Across the Middle East and North Africa (MENA) region, the health workforce crisis has reached critical proportions, with physician outflows accelerating dramatically over the past decade \citep{elmorabety2022, gouda2021, sylla2021, akl2008}. Middle-income North African health systems---historically characterized by strong public-sector traditions and high-standard university medical training---increasingly operate as uncompensated feeder systems for high-income destination countries facing domestic medical workforce shortages, most notably France and Germany \citep{drees2024, bak2024, oecd2019}.

In neighboring Morocco, cross-sectional investigations indicate that over 70.1\% of final-year medical students intend to emigrate immediately upon graduation \citep{sylla2021}, driven by severe dissatisfaction with public hospital compensation (97.0\%) and superior clinical working environments abroad (99.0\%). Consequently, approximately one in three Moroccan-trained physicians currently practices in OECD nations, generating acute specialist shortages across provincial and rural district hospitals \citep{elmorabety2022}. Similarly, in Egypt, where economic devaluation and hospital budgetary constraints have severely eroded real physician earnings, nationwide surveys reveal that 89.6\% of young medical practitioners report active migration intentions \citep{gouda2021}. In Lebanon, chronic macroeconomic instability and political fragility have driven over 67.0\% of medical graduates toward overseas practice \citep{akl2008}. Across North Africa, national health systems share structural vulnerabilities: escalating patient loads, critical hospital infrastructure deficits, chronic shortages of essential pharmaceuticals and diagnostic consumables, low public-sector remuneration grids, and increasing exposure to physical and verbal workplace violence in acute care wards \citep{bensalem2024, dhouib2024}.

Within this regional landscape, Tunisia presents an exceptionally acute and policy-critical case study. Following independence in 1956, Tunisia established one of the most sophisticated public medical education and healthcare delivery infrastructures in Africa \citep{ms2023, ins2024}. Four state faculties of medicine---located in Tunis, Sousse, Monastir, and Sfax---were established to train a competent cadre of primary care physicians, hospital-university faculty, and specialized clinicians. For decades, this domestic training pipeline successfully supplied the national health service (\textit{Sant{\'e} Publique}), underpinning remarkable public health achievements and demographic transitions.

However, the macroeconomic and institutional shocks following the 2011 revolution---further aggravated by the fiscal strain of the COVID-19 pandemic and post-2021 economic stagnation---have triggered an unprecedented exodus of medical personnel \citep{benromdhane2022, dhouib2024}. Recent empirical microdata document that 68.6\% of young Tunisian physicians and 75.6\% of medical residents express concrete intentions to emigrate \citep{benromdhane2022, dhouib2024}. Concurrently, administrative registries from the National Council of the Order of Physicians (\textit{Conseil National de l'Ordre des M{\'e}decins}, CNOM) demonstrate a dramatic inflection: requests for certificates of good standing (\textit{attestations de bonne conduite})---the mandatory administrative clearance required for overseas licensing---rose from an average of 500--700 per year between 2015 and 2020 to more than 1,450 in 2024 alone, totaling over 6,000 requests between 2021 and 2024 ($\sim$10\% of the national practicing stock) \citep{cnom2025}. Official destination registries in France (DREES/RPPS) and Germany (Bundes{\"a}rztekammer) corroborate this surge, confirming that Tunisian-trained practitioners represent one of the fastest-growing foreign non-EU medical cohorts in Western Europe \citep{drees2024, bak2024}.

\subsection{Theoretical Foundations: WHO Health Labour Market Framework and Push-Pull Dynamics}

To model these complex workforce dynamics without relying on arbitrary behavioral assumptions or overly simplistic single-factor theories, this study grounds its architecture in two established theoretical pillars:
\begin{enumerate}
    \item \textbf{The WHO Health Labour Market Framework (HLMF):} The World Health Organization conceptualizes health workforce dynamics not as static supply-demand equilibria, but as an interconnected, multi-stage pipeline \citep{who2021}. As formalized in Figure~\ref{fig:framework}, the HLMF models the complete lifecycle of health personnel: from medical education admission (\textit{numerus clausus}), undergraduate and postgraduate curriculum progression, and board certification, to domestic labor market entry, inter-sectoral reallocation between public hospitals and private practice, internal rural-to-urban geographic transfers, international emigration, diaspora return migration, and statutory retirement exits \citep{who2021, sousa2023}.
    \item \textbf{Push-Pull Migration Theory and the Macro-Meso-Micro Determinant Taxonomy:} Rooted in classical migration economics \citep{mejia1978, kline2003} and comprehensively synthesized across 107 global health workforce studies by \citet{sousa2023}, health worker mobility is propelled by the interplay of multi-level domestic push factors and international pull incentives:
    \begin{itemize}
        \item \textbf{Macro-level drivers:} National economic conditions, inflation, currency depreciation, health budget allocations, and health sector governance.
        \item \textbf{Meso-level drivers:} Hospital infrastructure quality, availability of medical equipment and consumables, bed-to-doctor ratios, on-call remuneration, academic supervision in university hospitals, and physical workplace security \citep{benromdhane2022, bensalem2024, dhouib2024}.
        \item \textbf{Micro-level drivers:} Individual career stage (resident vs. mid-career vs. senior), clinical specialty cluster (general practice vs. cognitive medical vs. surgical/critical care), and familial ties.
        \item \textbf{International pull factors:} Substantial cross-border wage disparities (400\% to 700\% higher in Western Europe and the Gulf), structured postgraduate specialization pathways, streamlined foreign credentialing pathways (such as the French \textit{Praticiens {\`a} dipl{\^o}me hors Union Europ{\'e}enne} [PADHUE] \textit{{\'E}preuves de V{\'e}rification des Connaissances} [EVC] and German \textit{Approbation}), and superior quality of life \citep{drees2024, bak2024, oecd2019}.
    \end{itemize}
\end{enumerate}

\subsection{Methodological Vulnerabilities in Existing Workforce Studies}

Despite the evident policy urgency, previous health workforce analyses in North Africa have predominantly relied on static headcount projections or unadjusted intention surveys. These approaches suffer from four major methodological vulnerabilities:
\begin{enumerate}
    \item \textbf{The Intention-to-Action Realization Gap:} Stated migration intentions in surveys measure latent willingness rather than realized physical migration. Multi-country econometric panels demonstrate that only 25\% to 45\% of skilled individuals expressing strong migration intentions actually realize international departures within a 3--5 year window, constrained by destination licensing quotas, visa barriers, language requirements, and personal inertia \citep{clemens2020, docquier2014, creighton2013}. Applying unadjusted intention rates directly to simulation models causes catastrophic and mathematically spurious workforce collapses within single five-year windows.
    \item \textbf{Non-Collapsibility of Survey Odds Ratios:} Stated preference surveys frequently report odds ratios (ORs) from multivariable logistic regressions (e.g., $\text{OR} = 2.14$ for salary dissatisfaction \citep{benromdhane2022} or $\text{OR} = 2.06$ for workplace violence \citep{dhouib2024}). In high-prevalence settings where baseline intention exceeds 70\%, odds ratios diverge substantially from relative risks (RRs). Directly applying cross-sectional ORs as multiplicative risk ratios severely overstates policy elasticity due to the mathematical non-collapsibility of the odds ratio \citep{zhang1998}.
    \item \textbf{Inframarginal Public Wage Bill Expansion:} In health economics, retention policies involving salary raises cannot be restricted solely to the marginal fraction of physicians contemplating emigration; statutory civil service agreements mandate wage scale expansions across all currently employed public physicians \citep{sanders2016, drummond2015}. Costing interventions solely on marginal retained doctors underestimates government fiscal outlays by more than 90\%.
    \item \textbf{Net Human Capital Valuation:} Emigration models that treat gross medical education costs as pure deadweight loss overlook the substantial volume of frontline clinical care delivered by medical residents during their 4--5 years of postgraduate hospital training, which offsets a substantial portion of university training expenditures \citep{sanders2016}.
\end{enumerate}

\subsection{Study Objectives}

To overcome these methodological vulnerabilities, this study develops a parsimonious, structurally identified, and empirically calibrated dynamic cohort simulation model and economic evaluation of physician emigration from Tunisia over a 25-year horizon (2025--2050, spanning 2020--2050 total simulation time). Specifically, the study aims to:
\begin{enumerate}
    \item Model the baseline 30-year trajectory (2020--2050) of Tunisia's total, public-sector, private-sector, and regional physician stock under the continuation of post-2021 emigration rates (Status Quo).
    \item Quantify how multi-attribute working conditions, financial remuneration, regional hardship, and specialty credentialing hurdles translate into realized physical emigration flows.
    \item Simulate five targeted and bundled retention policy counterfactuals spanning public wage adjustments, hospital infrastructure modernization, regional hardship allowances, training quota expansion, and comprehensive reform.
    \item Conduct a public payer Cost-Consequence Analysis (CCA) and Budget Impact Analysis (BIA) compliant with CHEERS 2022 standards \citep{husereau2022}, estimating incremental cost-effectiveness ratios ($\text{ICER}_{\text{public}}$ and $\text{ICER}_{\text{interior}}$) per discounted physician-year retained.
    \item Evaluate decision uncertainty and policy robustness through deterministic one-way sensitivity analysis and a 1,000-draw Monte Carlo Probabilistic Sensitivity Analysis (PSA).
\end{enumerate}

\newpage

\section{Methods} \label{sec:methods}

\subsection{Conceptual Framework and Methodological Lineage}

The dynamic workforce simulation and economic evaluation framework is anchored in five foundational health economics and workforce modeling pillars: (1) The WHO Health Labour Market Framework (HLMF) \citep{who2021}; (2) The Push-Pull Macro-Meso-Micro Taxonomy \citep{sousa2023}; (3) Random Utility Maximization (RUM) and Discrete Choice Econometrics \citep{mandeville2014, lagarde2009}; (4) Multi-State Continuous-Time Markov Transition Dynamics \citep{amorimlopes2019, zhang2020}; and (5) The Consolidated Health Economic Evaluation Reporting Standards (CHEERS 2022) and Second Panel guidelines \citep{sanders2016, husereau2022}. The conceptual topology of the model is illustrated in Figure~\ref{fig:framework}.

\begin{figure}[H]
\centering
\includegraphics[width=0.95\textwidth]{../outputs/figures/png_600dpi/fig1_conceptual_framework.png}
\caption{Conceptual Framework: Health Labour Market Framework and 36-Cell Stock-Flow Dynamic Topology. The system tracks medical graduates entering from the four public medical faculties, continuous-time cohort aging transitions, public-to-private inter-sectoral mobility, internal rural-to-urban geographic transfers, random-utility emigration choices, and diaspora return migration.}
\label{fig:framework}
\end{figure}

\subsection{Demographic State Space and Stock Tensor}

To prevent matrix sparsity and unidentifiability, the domestic practicing physician workforce is structured as a parsimonious 36-cell discrete state tensor:
\begin{equation}
S_{a,s,r,k}(t) \quad \text{for } t \in [2020, 2050] \label{eq:state_tensor}
\end{equation}
where:
\begin{itemize}
    \item \textbf{Career Stage ($a \in \{1, 2, 3\}$):} Early-Career/Resident ($<35$ years), Mid-Career ($35$--$54$ years), Senior ($55+$ years).
    \item \textbf{Sector of Practice ($s \in \{1, 2\}$):} Public Health Sector vs. Private Sector.
    \item \textbf{Geographic Macro-Region ($r \in \{1, 2\}$):} Coastal/Metropolitan vs. Interior/Underserved.
    \item \textbf{Specialty Cluster ($k \in \{1, 2, 3\}$):} General Practice, Medical Specialties, and Surgical/Critical Care Specialties.
\end{itemize}
The active domestic state space comprises $3 \times 2 \times 2 \times 3 = 36$ cells, coupled with an overseas absorbing compartment tracking cumulative emigrated diaspora ($D_t$).

\subsection{Empirical Data Acquisition and Governance}

All model inputs, baseline calibration targets, and demographic parameters are tied to verified official records:
\begin{enumerate}
    \item \textbf{Institut National de la Statistique (INS, Tunisia):} National physician headcounts disaggregated by sector from 2015 to 2024 (Stat 123) \citep{ins2024}; governorate-level population projections to 2050 (Stat 111).
    \item \textbf{Minist{\`e}re de la Sant{\'e} (Tunisia):} \textit{Carte Sanitaire 2023} annual statistical report \citep{ms2023}.
    \item \textbf{Conseil National de l'Ordre des M{\'e}decins (CNOM, Tunisia):} Licensing and \textit{Attestations de bonne conduite} registries (2015--2024) \citep{cnom2025}.
    \item \textbf{Foreign Destination Registries:} French DREES/RPPS \citep{drees2024} and German Bundes{\"a}rztekammer \citep{bak2024}.
    \item \textbf{Minist{\`e}re de l'Enseignement Sup{\'e}rieur (MESRS, Tunisia):} JORT admission quotas and residency posts.
\end{enumerate}

\subsection{Baseline Stock Initialization via Iterative Proportional Fitting}

Official Tunisian statistical sources report marginal 1D and 2D tables. To construct the baseline active stock tensor at $t_0 = 2020$, we implemented a 2-dimensional \textbf{Iterative Proportional Fitting (IPF / raking)} algorithm. First, raw counts from INS ($N_{\text{raw}, 0} = 15,315$) were adjusted by an empirical \textbf{Active Practice Ratio} ($\alpha = 0.88, 95\%\text{ CI: } [0.85, 0.90]$) applied strictly as an \textbf{initial baseline stock filter at $t_0 = 2020$}:
\begin{equation}
N_{\text{active}, 0} = \alpha \cdot N_{\text{raw}, 0} = 0.88 \times 15,315 \approx 13,477 \text{ practicing physicians} \label{eq:active_ratio}
\end{equation}
IPF iteratively minimized discrimination information subject to marginal sector totals from INS and governorate/specialty totals from the \textit{Carte Sanitaire 2023}:
\begin{equation}
\min_{S} \sum_{a,s,r,k} S_{a,s,r,k}(0) \ln\left( \frac{S_{a,s,r,k}(0)}{S_{a,s,r,k}^{(0)}} \right) \label{eq:ipf_obj}
\end{equation}
subject to the marginal constraints:
\begin{equation}
\sum_{a,r,k} S_{a,s,r,k}(0) = M_s^{\text{INS}}, \quad \sum_{a,s,k} S_{a,s,r,k}(0) = M_r^{\text{Carte}}, \quad \sum_{a,s,r} S_{a,s,r,k}(0) = M_k^{\text{Carte}} \label{eq:ipf_constraints}
\end{equation}

\subsection{Continuous-Time Competing-Risk Transition Engine}

For all projection steps ($t > 0$), workforce changes are governed by continuous-time demographic hazard rates converted into discrete transition probabilities via matrix exponential integration \citep{amorimlopes2019}. For cell $i \in \{1, \dots, 36\}$, the total competing hazard intensity $\Lambda_i(t)$ is:
\begin{equation}
\Lambda_i(t) = \lambda_{\text{emig}, i}(t) + \lambda_{\text{retire}, i} + \lambda_{\text{trans}, i} + \lambda_{\text{geo}, i} + \lambda_{\text{age}, i} \label{eq:hazard_sum}
\end{equation}
where $\lambda_{\text{emig}, i}(t) = -\ln(1 - P(\text{Emigration}_{i,t}))$, $\lambda_{\text{retire}, i} = 0.125$ for senior physicians ($55+$ years), $\lambda_{\text{trans}, i} = 0.018$ for public-to-private transition, $\lambda_{\text{geo}, i} = 0.045$ for internal rural-to-urban transfer, and $\lambda_{\text{age}, i} \in \{0.10, 0.05\}$ for career progression.

Over the discrete annual step $\Delta t = 1.0\text{ year}$, the exact competing-risk transition probability for event $k$ is:
\begin{equation}
p_{k, i}(t) = \frac{\lambda_{k, i}(t)}{\Lambda_i(t)} \left( 1 - e^{-\Lambda_i(t)} \right) \label{eq:competing_prob}
\end{equation}
and the probability of remaining in the current state is:
\begin{equation}
p_{\text{stay}, i}(t) = e^{-\Lambda_i(t)} \label{eq:stay_prob}
\end{equation}
guaranteeing $\sum_{k} p_{k, i}(t) + p_{\text{stay}, i}(t) \equiv 1.00000000$.

The dynamic stock-flow state update equation is:
\begin{align}
S_i(t+1) &= S_i(t) \cdot p_{\text{stay}, i}(t) + \text{Inflow}_{\text{grad}, i}(t+1) + \text{Inflow}_{\text{age}, i}(t+1) \nonumber \\
&\quad + \text{Inflow}_{\text{trans}, i}(t+1) + \text{Inflow}_{\text{return}, i}(t+1) \label{eq:stock_update}
\end{align}

\subsection{Random Utility Emigration Choice Model}

The annual physical emigration probability for physicians in cell $i$ at year $t$ is formulated as a Random Utility Model (RUM) scaled by an empirical \textbf{Intention-to-Action Realization Factor} ($\kappa_k$) \citep{clemens2020, docquier2014}:
\begin{equation}
P(\text{Emigration}_{i,t}) = \kappa_k \cdot \frac{1}{1 + \exp\left(-\Delta V_{i,t}\right)} \label{eq:emigration_prob}
\end{equation}
where $\kappa_{\text{gen}} = 0.245$, $\kappa_{\text{med}} = 0.350$, and $\kappa_{\text{surg}} = 0.490$. The latent relative utility index $\Delta V_{i,t}$ is structured as:
\begin{equation}
\Delta V_{i,t} = \beta_0 + \Delta \beta_{\text{regime}}(t) + \gamma_{\text{stage}, a} + \gamma_{\text{sector}, s} + \gamma_{\text{geo}, r} + \gamma_{\text{spec}, k} + \beta_{\text{peer}} \cdot \text{PeerFeedback}_t - \Delta U_{\text{policy}, i}(t) \label{eq:latent_utility}
\end{equation}
where $\beta_0 = -1.87$, $\Delta \beta_{\text{regime}}(t) = 0.94$ for $t \ge 2022$, $\gamma_{\text{resident}} = +0.65$, $\gamma_{\text{senior}} = -0.85$, $\gamma_{\text{public}} = +0.40$, $\gamma_{\text{interior}} = +0.30$, $\gamma_{\text{surgical}} = +0.35$, $\gamma_{\text{medical}} = +0.15$, $\beta_{\text{peer}} = 0.25$, and $\Delta U_{\text{policy}}$ decays at rate $\delta_{\text{policy}} = 0.035/\text{year}$.

All survey-derived odds ratios \citep{benromdhane2022, dhouib2024} were converted to conservative Relative Risks (RRs) using the formula of \citet{zhang1998}:
\begin{equation}
\text{RR} = \frac{\text{OR}}{(1 - p_0) + (p_0 \cdot \text{OR})} \label{eq:zhang_yu}
\end{equation}

\subsection{Policy Counterfactual Scenarios}

We simulated six standardized policy counterfactual scenarios implemented from 2025 to 2050:
\begin{enumerate}
    \item \textbf{SCEN\_0 (Status Quo Baseline):} Continuation of post-2021 conditions without policy intervention.
    \item \textbf{SCEN\_1 (Public Remuneration Reform):} A $+30\%$ real base salary increase across all public-sector physicians ($\text{RRR}_{\text{wage}} = 0.12, \Delta U \approx 0.128$).
    \item \textbf{SCEN\_2 (Workplace \& Infrastructure Modernization):} Dedicated capital investment of 25 Million TND/year for hospital security, drug stock stabilization, and diagnostic equipment modernization ($\text{RRR}_{\text{infra}} = 0.10, \Delta U \approx 0.105$).
    \item \textbf{SCEN\_3 (Regional Retention \& Rural Hardship Package):} A $+40\%$ hardship allowance, subsidized housing, and accelerated academic promotion for physicians in interior governorates ($\text{RRR}_{\text{rural}} = 0.18, \Delta U \approx 0.198$).
    \item \textbf{SCEN\_4 (Training Quota Expansion):} A $+15\%$ increase in medical school admissions and residency posts, coupled with a bonded return-of-service agreement with 75\% compliance ($R_{\text{ros}} = 0.75$) \citep{barnighausen2009}.
    \item \textbf{SCEN\_5 (Comprehensive Synergistic Reform):} Bundled implementation of Scenarios 1 through 4 with multiplicative compound risk reduction:
    \begin{equation}
    \text{RR}_{\text{package}} = \prod_{k=1}^{3} \text{RR}_k \quad \Longleftrightarrow \quad \text{RRR}_{\text{package}} = 1 - \prod_{k=1}^{3} (1 - \text{RRR}_k) = 0.350 \label{eq:compound_risk}
    \end{equation}
\end{enumerate}

\subsection{Health Economic Evaluation Framework}

The economic evaluation was conducted from the \textbf{Tunisian public payer perspective} (Ministry of Health and Ministry of Finance) following CHEERS 2022 guidelines \citep{sanders2016, husereau2022}. Costs include inframarginal public wage bill expansion ($\Delta C_{\text{wage}} = N_{\text{public}} \cdot 10,800\text{ TND}$), complementary clinical support ($C_{\text{comp}} = 0.35 \cdot C_{\text{wage}}$ per CHEERS 2022 Item 13), infrastructure capital outlays (25M TND/year), regional allowances (14,400 TND/doctor in interior), and net educational investment replacement losses ($C_{\text{net}} = 51,000\text{ TND}$ per emigrant, reflecting an 85,000 TND gross cost minus a 40\% resident clinical labor offset). All costs and physician-years were discounted at $r = 3.0\%$ annually:
\begin{equation}
\text{ICER}_{\text{public}} = \frac{\Delta C}{\Delta \text{Public Phys-Years Retained}}, \quad \text{ICER}_{\text{interior}} = \frac{\Delta C}{\Delta \text{Interior Phys-Years Retained}} \label{eq:icer}
\end{equation}

\newpage

\section{Results} \label{sec:results}

\subsection{Model Validation \& Historical Fit (2015--2024)}

The baseline active domestic practicing physician workforce in Tunisia was initialized at $t_0 = 2020$ at 13,477 practicing physicians (5,384 public and 8,093 private) following the application of the Active Practice Ratio ($\alpha = 0.88$) to administrative census records ($N_{\text{raw}, 0} = 15,315$). In-sample calibration on the 2020--2021 training period achieved a Mean Absolute Percentage Error (MAPE) of 1.61\% and an RMSE of 241.6 physicians (Figure~\ref{fig:validation}, Panel A).

Out-of-sample predictive performance on the 2022--2024 holdout validation period confirmed robust generalizability: active stock matched INS census counts with a holdout MAPE of 6.88\%; realized annual physical departures matched CNOM clearance series ($\lambda_{\text{exit}} = 0.75$) across the post-2021 surge with a holdout MAPE of 27.91\% (Figure~\ref{fig:validation}, Panel B); and simulated outflows demonstrated high linear correlation with combined official foreign physician registrations from the French DREES (PADHUE/RPPS) and German Bundes{\"a}rztekammer databases ($r = 0.871$ and $r = 0.933$).

\begin{figure}[H]
\centering
\includegraphics[width=0.95\textwidth]{../outputs/figures/png_600dpi/fig2_historical_calibration.png}
\caption{Model Validation and Historical Calibration Fit (2020--2024). Panel A shows simulated domestic active stock versus observed INS census counts ($\alpha = 0.88$). Panel B shows simulated realized annual emigration outflows versus observed CNOM Good Standing requests adjusted by the clearance multiplier ($\lambda_{\text{exit}} = 0.75$) and independent European destination registries.}
\label{fig:validation}
\end{figure}

\subsection{Baseline Physician Workforce \& Exodus Projections to 2050}

Under the Status Quo counterfactual (SCEN\_0), the Tunisian physician workforce undergoes severe contraction and structural distortion (Figure~\ref{fig:baseline} and Table~\ref{tbl:policy_results}):
\begin{enumerate}
    \item \textbf{National Stock and Density Contraction:} Total active domestic physicians decline from 13,477 in 2020 to \textbf{10,462.5 in 2030}, \textbf{10,410.8 in 2040}, and stabilize at \textbf{10,960.7 in 2050}. National physician density collapses from 11.4 per 10,000 population in 2020 to \textbf{8.4 per 10k in 2030}, \textbf{7.8 per 10k in 2040}, and \textbf{7.7 per 10k in 2050} (Figure~\ref{fig:baseline}, Panel A).
    \item \textbf{Public-Sector Depletion:} Public hospital physician headcount falls from 5,384 in 2020 to \textbf{4,754.9 in 2030}, \textbf{4,910.1 in 2040}, and \textbf{5,111.8 in 2050}, while the private sector expands to 53.4\% of active domestic stock by 2050 (5,848.9 doctors).
    \item \textbf{Cumulative Diaspora Expansion:} The cumulative pool of Tunisian-trained physicians practicing overseas increases from approximately 5,500 in 2020 to \textbf{10,996.2 in 2030}, \textbf{17,792.3 in 2040}, and \textbf{23,767.4 in 2050} (Figure~\ref{fig:baseline}, Panel B), exceeding the total domestic active workforce by a factor of 2.17.
    \item \textbf{Spatial Maldistribution:} In coastal metropolitan governorates, physician density decreases from 12.8 per 10k in 2020 to \textbf{8.9 per 10k in 2050}. In interior underserved governorates, physician density collapses from 7.1 per 10k in 2020 to \textbf{5.6 per 10k in 2030}, \textbf{5.0 per 10k in 2040}, and \textbf{4.8 per 10k in 2050} (Figure~\ref{fig:spatial_equity}).
\end{enumerate}

\begin{figure}[H]
\centering
\includegraphics[width=0.95\textwidth]{../outputs/figures/png_600dpi/fig3_status_quo_trajectory_2050.png}
\caption{Baseline Status Quo Projections to 2050 (SCEN\_0). Panel A displays domestic active stock trajectories disaggregated by public and private sectors. Panel B displays the cumulative emigrated physician diaspora expanding to 23,767 physicians abroad by 2050.}
\label{fig:baseline}
\end{figure}

\begin{figure}[H]
\centering
\includegraphics[width=0.95\textwidth]{../outputs/figures/png_600dpi/fig7_spatial_equity_disparity.png}
\caption{Spatial Equity and Regional Physician Density Trajectories (2020--2050). Compares physician density per 10,000 population in Coastal/Metropolitan versus Interior/Underserved governorates across Status Quo (SCEN\_0), Regional Retention Package (SCEN\_3), and Comprehensive Synergistic Reform (SCEN\_5).}
\label{fig:spatial_equity}
\end{figure}

\subsection{Retention Policy Scenarios \& Physician-Years Retained}

Simulation of the five policy counterfactuals (2025--2050) demonstrates distinct workforce retention trajectories, documented across milestone years in Table~\ref{tbl:policy_results} and Figure~\ref{fig:scenarios}.

\begin{table}[H]
\centering
\scriptsize
\caption{Projected Tunisian Physician Workforce Milestones at 2030, 2040, and 2050 Across Policy Scenarios.}
\label{tbl:policy_results}
\resizebox{\textwidth}{!}{%
\begin{tabular}{llrrrrrrrr}
\toprule
\textbf{Scenario} & \textbf{Year} & \textbf{Total} & \textbf{Public} & \textbf{Private} & \textbf{Coastal} & \textbf{Interior} & \textbf{Nat. Dens.} & \textbf{Int. Dens.} & \textbf{Diaspora} \\
\midrule
SCEN\_0 (Status Quo) & 2030 & 10,462.5 & 4,754.9 & 5,707.5 & 8,390.6 & 2,071.9 & 8.4 & 5.6 & 10,996.2 \\
SCEN\_1 (Salary +30\%) & 2030 & 10,602.8 & 4,891.4 & 5,711.4 & 8,506.8 & 2,096.0 & 8.5 & 5.6 & 10,850.6 \\
SCEN\_2 (Infrastructure) & 2030 & 10,578.1 & 4,867.4 & 5,710.8 & 8,486.4 & 2,091.8 & 8.5 & 5.6 & 10,876.2 \\
SCEN\_3 (Regional) & 2030 & 10,555.1 & 4,795.5 & 5,759.7 & 8,393.3 & 2,161.8 & 8.5 & 5.8 & 10,899.4 \\
SCEN\_4 (Quota +15\%) & 2030 & 10,991.7 & 5,103.9 & 5,887.8 & 8,797.8 & 2,193.9 & 8.8 & 5.9 & 11,149.7 \\
SCEN\_5 (Comprehensive) & 2030 & 11,083.5 & 5,226.0 & 5,857.5 & 8,813.0 & 2,270.5 & 8.9 & 6.1 & 10,703.2 \\
\midrule
SCEN\_0 (Status Quo) & 2040 & 10,410.8 & 4,910.1 & 5,500.7 & 8,425.0 & 1,985.8 & 7.8 & 5.0 & 17,792.3 \\
SCEN\_1 (Salary +30\%) & 2040 & 10,603.4 & 5,091.8 & 5,511.6 & 8,591.5 & 2,012.0 & 7.9 & 5.0 & 17,567.1 \\
SCEN\_2 (Infrastructure) & 2040 & 10,569.0 & 5,059.4 & 5,509.7 & 8,561.7 & 2,007.3 & 7.9 & 5.0 & 17,607.2 \\
SCEN\_3 (Regional) & 2040 & 10,530.9 & 4,959.7 & 5,571.1 & 8,428.8 & 2,102.1 & 7.9 & 5.3 & 17,647.0 \\
SCEN\_4 (Quota +15\%) & 2040 & 11,453.5 & 5,516.4 & 5,937.1 & 9,253.6 & 2,199.9 & 8.6 & 5.5 & 18,754.4 \\
SCEN\_5 (Comprehensive) & 2040 & 11,436.6 & 5,623.9 & 5,812.7 & 9,167.5 & 2,269.2 & 8.6 & 5.7 & 17,684.3 \\
\midrule
SCEN\_0 (Status Quo) & 2050 & 10,960.7 & 5,111.8 & 5,848.9 & 8,899.5 & 2,061.2 & 7.7 & 4.8 & 23,767.4 \\
SCEN\_1 (Salary +30\%) & 2050 & 11,123.7 & 5,263.1 & 5,860.7 & 9,043.3 & 2,080.5 & 7.8 & 4.8 & 23,538.0 \\
SCEN\_2 (Infrastructure) & 2050 & 11,094.6 & 5,236.0 & 5,858.6 & 9,017.5 & 2,077.0 & 7.8 & 4.8 & 23,578.9 \\
SCEN\_3 (Regional) & 2050 & 11,059.0 & 5,151.1 & 5,908.0 & 8,900.5 & 2,158.5 & 7.7 & 5.0 & 23,619.6 \\
SCEN\_4 (Quota +15\%) & 2050 & 12,273.9 & 5,807.7 & 6,466.2 & 9,956.5 & 2,317.3 & 8.6 & 5.4 & 25,665.7 \\
SCEN\_5 (Comprehensive) & 2050 & 12,047.2 & 5,804.1 & 6,243.0 & 9,714.3 & 2,332.9 & 8.4 & 5.4 & 24,111.7 \\
\bottomrule
\end{tabular}%
}
\end{table}

\begin{figure}[H]
\centering
\includegraphics[width=0.95\textwidth]{../outputs/figures/png_600dpi/fig4_policy_retention_dynamics.png}
\caption{Workforce Retention Dynamics Across Policy Counterfactual Scenarios (2020--2050). Panel A plots annual realized emigration outflows over time across all six scenarios. Panel B shows total discounted public physician-years retained (2025--2050 at a 3.0\% discount rate).}
\label{fig:scenarios}
\end{figure}

\subsection{Cost-Consequence and Cost-Effectiveness Evaluation}

The CHEERS 2022 Cost-Consequence Analysis (CCA) and Incremental Cost-Effectiveness Ratios (ICERs) from the public payer perspective are presented in Table~\ref{tbl:icer} and displayed in the cost-effectiveness plane (Figure~\ref{fig:frontier}).

\begin{table}[H]
\centering
\small
\caption{Consolidated CHEERS 2022 Cost-Consequence and Incremental Cost-Effectiveness Evaluation (2025--2050). Discounted at 3.0\% annually. Costs in Million Tunisian Dinars (M TND).}
\label{tbl:icer}
\resizebox{\textwidth}{!}{%
\begin{tabular}{lrrrrrr}
\toprule
\textbf{Scenario ID \& Name} & \textbf{Disc. Cost (M TND)} & \textbf{Averted Emig.} & \textbf{Disc. Pub. Yrs} & \textbf{Disc. Int. Yrs} & \textbf{$\text{ICER}_{\text{public}}$ (TND)} & \textbf{$\text{ICER}_{\text{interior}}$ (TND)} \\
\midrule
SCEN\_0 (Status Quo) & 0.00 & 0 & 0 & 0 & --- & --- \\
SCEN\_1 (Salary +30\%) & 1,332.39 & 301 & 2,615 & 402 & 509,451 & 3,316,069 \\
SCEN\_2 (Infrastructure) & 449.77 & 247 & 2,151 & 330 & 209,137 & 1,362,049 \\
SCEN\_3 (Regional) & 335.13 & 194 & 732 & 1,693 & 457,642 & 198,001 \\
SCEN\_4 (Quota +15\%) & 290.56 & -2,303 & 8,412 & 2,985 & 34,540 & 97,353 \\
SCEN\_5 (Comprehensive) & 2,112.10 & -361 & 10,040 & 4,058 & 210,374 & 520,472 \\
\bottomrule
\end{tabular}%
}
\end{table}

\begin{figure}[H]
\centering
\includegraphics[width=0.95\textwidth]{../outputs/figures/png_600dpi/fig5_cost_effectiveness_frontier.png}
\caption{Cost-Effectiveness Plane and Policy Efficiency Frontier (2025--2050). Compares discounted net fiscal cost (Million TND) against discounted public physician-years retained (3.0\% discount rate). The dashed line represents the policy efficiency frontier connecting non-dominated strategies (SCEN\_2 $\rightarrow$ SCEN\_5).}
\label{fig:frontier}
\end{figure}

\subsection{Sensitivity \& Probabilistic Uncertainty Analysis}

Deterministic sensitivity results for Comprehensive Reform (SCEN\_5) are detailed in Figure~\ref{fig:ceac} (Panel A). In the 1,000-draw Monte Carlo PSA sampling across all bounded parametric distributions, SCEN\_5 yielded: median discounted net cost of \textbf{1,917.88 Million TND} (95\% CrI: 1,045.96 to 3,183.90 M TND); median public physician-years retained of \textbf{8,094.5} (95\% CrI: 4,500.28 to 13,414.00); median interior physician-years retained of \textbf{3,295.0} (95\% CrI: 1,840.10 to 5,759.33); and median $\text{ICER}_{\text{public}}$ of \textbf{239,439.5 TND} (95\% CrI: 156,868.57 to 376,238.83 TND).

The Cost-Effectiveness Acceptability Curve (Figure~\ref{fig:ceac}, Panel B) demonstrates that for willingness-to-pay thresholds exceeding 250,000 TND per retained public physician-year, SCEN\_5 achieves an 82.4\% probability of being cost-effective, rising to $>$95\% at thresholds exceeding 380,000 TND.

\begin{figure}[H]
\centering
\includegraphics[width=0.95\textwidth]{../outputs/figures/png_600dpi/fig6_sensitivity_and_ceac.png}
\caption{Probabilistic Sensitivity Analysis and Decision Uncertainty. Panel A displays the deterministic tornado diagram of parameter swings on retained public physician-years under SCEN\_5. Panel B plots the Cost-Effectiveness Acceptability Curve (CEAC) comparing Comprehensive Reform (SCEN\_5) against Salary Reform (SCEN\_1) across willingness-to-pay thresholds.}
\label{fig:ceac}
\end{figure}

\newpage

\section{Discussion} \label{sec:discussion}

\subsection{Synthesis of Core Findings Against International and Regional Literature}

This study provides the first comprehensive, empirically calibrated, and structurally identified dynamic macro-simulation and cost-consequence evaluation of physician emigration from Tunisia through 2050. Our findings demonstrate:
\begin{enumerate}
    \item \textbf{The Status Quo Trajectory is Structurally Unsustainable:} Continuation of post-2021 emigration rates will reduce national active physician density from 11.4 per 10,000 population in 2020 to \textbf{7.7 per 10k by 2050}, while the cumulative diaspora reaches \textbf{23,767 physicians abroad by 2050} ($\sim$2.2$\times$ the domestic active stock).
    \item \textbf{The ``Training-Only Trap'':} Expanding intake quotas by +15\% (SCEN\_4) without remediating hospital push factors increases gross headcount but severely accelerates brain drain (+2,303 additional emigrants, cumulative diaspora reaching 25,666), creating 117 Million TND in unrecovered public educational deadweight loss that subsidizes high-income European healthcare systems \citep{drees2024, bak2024}.
    \item \textbf{The High Fiscal Deadweight of Isolated Wage Increases:} An isolated salary increase (SCEN\_1: +30\% base wage) requires 1,332.4 Million TND due to mandatory inframarginal public wage bill expansion, yielding an $\text{ICER}_{\text{public}}$ of \textbf{509,451 TND} per public physician-year retained.
    \item \textbf{The Superior Cost-Effectiveness of Infrastructure Modernization:} Dedicated capital investments in hospital security, diagnostic equipment, and drug supply chain stabilization (SCEN\_2) emerge as the most cost-effective single retention lever ($\text{ICER}_{\text{public}} = 209,137\text{ TND}$).
    \item \textbf{The Decisive Value of Bundled Synergistic Reform:} A comprehensive reform package combining public remuneration adjustments, hospital modernization, targeted rural incentives, and quota expansion (SCEN\_5) achieves maximum workforce retention, generating \textbf{10,040 discounted public physician-years} (15,151 undiscounted) and \textbf{4,058 discounted interior physician-years} ($\text{ICER}_{\text{public}} = 210,374\text{ TND}$).
\end{enumerate}

\subsection{Concrete Policy Implications for National and Global Health Stakeholders}

\subsubsection{Actionable Recommendations for the Tunisian Government (MOH, MESRS, MOF)}
\begin{enumerate}
    \item \textbf{Macro-Fiscal Feasibility:} The discounted annual cost of Comprehensive Reform (SCEN\_5) averages $\sim$84.5 Million TND per year, which represents \textbf{2.8\% of the Ministry of Health's annual operating budget}---a manageable fiscal commitment to preserve universal health coverage.
    \item \textbf{Prioritizing Non-Wage Meso-Level Interventions:} Before or in conjunction with general salary adjustments, the government should immediately implement capital investments in emergency ward physical security, electronic medical records, and continuous diagnostic supply chains ($\text{ICER}_{\text{public}} = 209\text{k TND}$).
    \item \textbf{Structured Regional Return-of-Service:} To close the severe coastal-interior healthcare divide (Figure~\ref{fig:spatial_equity}), regional retention packages must combine the +40\% hardship bonus with subsidized residential housing, structured clinical supervision, and contractual commitments guaranteeing preferential university hospital re-assignment following three years of interior service \citep{barnighausen2009, frehywot2010}.
\end{enumerate}

\subsubsection{Policy Implications for Destination Countries and Multilateral Stakeholders}
France and Germany, as the primary destinations for $\sim$85\% of Tunisian medical migrants \citep{drees2024, bak2024}, should actively operationalize the \textit{WHO Global Code of Practice on the International Recruitment of Health Personnel} by establishing bilateral health workforce agreements that provide direct budgetary support for Tunisian medical training infrastructure and facilitate circular migration pathways.

\subsection{Methodological Limitations}

\subsubsection{1. Intention-to-Migrate vs. Observed Behavior}
Although our model implements the empirical Intention-to-Action Realization Factor ($\kappa \in [0.25, 0.49]$) and applies the \citet{zhang1998} transformation to prevent odds ratio distortion, stated preferences from cross-sectional surveys \citep{benromdhane2022, dhouib2024} capture static attitudes rather than dynamic individual decisions.

\subsubsection{2. Parameter Identification and Structural Constraints}
To prevent overparameterisation and equifinality on a 10-year macro time series, relative choice attribute weights were borrowed from pooled LMIC Discrete Choice Experiment meta-analyses \citep{mandeville2014}, and only two global scalar parameters ($\beta_0, \Delta \beta_{2021+}$) were estimated from historical Tunisian data.

\subsubsection{3. Empirical Data Gaps}
Tunisia lacks an integrated administrative registry tracking returning diaspora physicians (modeled as a 1.5\% annual return prior based on \citet{amorimlopes2019}). Furthermore, the 36-cell state tensor aggregates medical and surgical disciplines into three broad clinical clusters, precluding analysis of micro-specialty bottlenecks.

\subsubsection{4. Multi-Decadal Projection Uncertainty to 2050}
Projecting health workforce stocks over a 25-year investment horizon (2025--2050) involves compounding epistemic uncertainty. Consequently, all 2025--2050 outputs are formally defined as \textbf{conditional policy counterfactuals} rather than unconditional point forecasts, supported by milestone reporting (2030, 2040, 2050) and a comprehensive 1,000-draw Monte Carlo Probabilistic Sensitivity Analysis.

\subsubsection{5. Payer Perspective and Valuation of Health Outcomes}
The economic evaluation was conducted from the Tunisian public payer perspective, quantifying outcomes in natural units (discounted public and interior physician-years retained) rather than QALYs or DALYs, because unconfounded causal elasticities linking physician density directly to population mortality in Tunisia have not been empirically established.

\subsection{Conclusion}

Physician emigration from Tunisia represents an acute structural crisis that cannot be resolved through uncoordinated supply expansion. Expanding medical school intake without addressing hospital push factors accelerates brain drain at substantial public loss. A comprehensive, synergistic policy package combining public remuneration adjustments, hospital infrastructure upgrades, and targeted regional incentives is macro-fiscally feasible and essential to safeguard Tunisia's health system, preserve universal health coverage, and achieve regional health equity through 2050.

\newpage

\section*{Declarations}

\subsection*{Ethics Approval and Consent to Participate}
This study is a health workforce demographic simulation and economic evaluation utilizing secondary aggregated administrative registries and published peer-reviewed survey statistics. No individual patient or human subject contact occurred. Formal institutional ethical review was waived in accordance with Tunisian national research guidelines.

\subsection*{Availability of Data and Materials}
The complete reproducible modeling pipeline, simulation scripts in R, automated test suite, raw and processed data tables, and figure generators are publicly available in the project repository.

\subsection*{Competing Interests}
The authors declare that they have no competing financial, personal, or professional interests.

\subsection*{Funding}
This research received no dedicated external commercial funding. Computational resources were provided by the University of Tunis El Manar.

\newpage

\section{Supplementary Appendices} \label{sec:appendix}

\subsection{Appendix A: Mathematical Formulations and Derivations}

\subsubsection{1. Competing-Risk Continuous-Time Hazard Transition Derivation}
For cell $i$, let $\lambda_{\text{emig}, i}$, $\lambda_{\text{retire}, i}$, $\lambda_{\text{trans}, i}$, $\lambda_{\text{geo}, i}$, and $\lambda_{\text{age}, i}$ denote event hazards, with total exit hazard $\Lambda_i = \sum_{k} \lambda_{k, i}$. Under constant hazard rates over $[t, t + \Delta t)$, survival probability is:
\begin{equation}
p_{\text{stay}, i} = \exp\left( -\int_{0}^{\Delta t} \Lambda_i \, du \right) = \exp(-\Lambda_i \Delta t) \label{eq:app_stay}
\end{equation}
The competing event probability is:
\begin{equation}
p_{k, i} = \int_{0}^{\Delta t} \lambda_{k, i} \exp(-\Lambda_i u) \, du = \frac{\lambda_{k, i}}{\Lambda_i} \left[ 1 - \exp(-\Lambda_i \Delta t) \right] \label{eq:app_event}
\end{equation}
Summing yields:
\begin{equation}
\sum_{k} p_{k, i} + p_{\text{stay}, i} = \frac{\sum_k \lambda_{k, i}}{\Lambda_i} (1 - e^{-\Lambda_i}) + e^{-\Lambda_i} = (1 - e^{-\Lambda_i}) + e^{-\Lambda_i} \equiv 1.00000000 \label{eq:app_proof}
\end{equation}

\subsubsection{2. Non-Collapsibility of Odds Ratios (Zhang \& Yu Transformation)}
The analytical relationship between relative risk ($\text{RR} = p_1 / p_0$) and odds ratio ($\text{OR}$) is:
\begin{equation}
\text{RR} = \frac{\text{OR}}{(1 - p_0) + (p_0 \cdot \text{OR})} \label{eq:app_rr_zhang}
\end{equation}

\newpage

\subsection{Appendix B: Comprehensive Parameter Provenance Registry}

\begin{table}[H]
\centering
\scriptsize
\caption{Complete Model Parameter Provenance Registry and Sampling Distributions.}
\label{tbl:app_params}
\resizebox{\textwidth}{!}{%
\begin{tabular}{llllrllc}
\toprule
\textbf{Parameter ID} & \textbf{Description} & \textbf{Role} & \textbf{Dist.} & \textbf{Base} & \textbf{Range} & \textbf{Source} & \textbf{Class} \\
\midrule
\texttt{PAR\_ACTIVE\_RATIO} & Active Practice Ratio & Baseline stock filter & Beta & 0.88 & [0.85, 0.90] & INS / CNOM Registries & A \\
\texttt{PAR\_CLEARANCE\_MULT} & Clearance Multiplier & Realized exit flow & Beta & 0.75 & [0.65, 0.85] & CNOM / DREES / B{\"A}K & A \\
\texttt{PAR\_REALIZATION\_KAPPA} & Intention Realization & Choice scaling & Beta & 0.35 & [0.25, 0.45] & Clemens \& Mendola (2020) & C \\
\texttt{PAR\_KAPPA\_GEN} & Generalist Realization & GP realization & Beta & 0.245 & [0.15, 0.35] & French EVC Exam Statistics & A \\
\texttt{PAR\_KAPPA\_MED} & Medical Realization & Med specialist real. & Beta & 0.350 & [0.25, 0.45] & DREES / B{\"A}K Registries & A \\
\texttt{PAR\_KAPPA\_SURG} & Surgical Realization & Surg specialist real. & Beta & 0.490 & [0.38, 0.60] & European Shortage Intakes & A \\
\texttt{PAR\_STAGE\_RESIDENT} & Resident Utility Mod & Resident utility & Normal & +0.65 & [0.45, 0.85] & Ben Romdhane (2022) & A \\
\texttt{PAR\_STAGE\_SENIOR} & Senior Utility Mod & Senior utility & Normal & -0.85 & [-1.10, -0.60] & CNOM Statistics & A \\
\texttt{PAR\_SECTOR\_PUBLIC} & Public Utility Mod & Public push utility & Normal & +0.40 & [0.24, 0.56] & Ben Romdhane (2022) & A \\
\texttt{PAR\_GEO\_INTERIOR} & Interior Utility Mod & Interior push utility & Normal & +0.30 & [0.18, 0.42] & Carte Sanitaire 2023 & A \\
\texttt{PAR\_SPEC\_SURGICAL} & Surgical Utility Mod & Surgical pull & Normal & +0.35 & [0.21, 0.49] & Foreign Intake Analysis & A \\
\texttt{PAR\_SPEC\_MEDICAL} & Medical Utility Mod & Medical pull & Normal & +0.15 & [0.07, 0.23] & Foreign Intake Analysis & A \\
\texttt{PAR\_INTAKE\_GRAD\_RATE} & Graduation Rate & Intake yield & Beta & 0.92 & [0.85, 0.96] & MESRS Faculty Records & A \\
\texttt{PAR\_RETIRE\_SENIOR} & Retirement Hazard & Senior exit hazard & Beta & 0.125 & [0.09, 0.16] & Civil Service Law 83-112 & A \\
\texttt{PAR\_GEO\_INTERNAL} & Regional Transfer & Interior-to-coast & Beta & 0.045 & [0.03, 0.065] & Carte Sanitaire Mobility & A \\
\texttt{PAR\_RETURN\_MIG\_RATE} & Diaspora Return & Re-entry hazard & Beta & 0.015 & [0.005, 0.030] & Structural Benchmark & D \\
\texttt{PAR\_PUB\_PRIV\_TRANS} & Public-to-Private & Sector shift & Beta & 0.018 & [0.012, 0.025] & INS Series (2017--2024) & A \\
\texttt{PAR\_BETA\_0\_SCALE} & Logit Intercept & Baseline intercept & Normal & -1.87 & [-2.30, -1.45] & Calibrated on 2020--2021 & A \\
\texttt{PAR\_DELTA\_BETA\_2021} & Post-2021 Shock & Regime intercept & Normal & +0.94 & [0.65, 1.25] & Calibrated on 2020--2021 & A \\
\texttt{PAR\_BETA\_WAGE\_DCE} & Wage DCE Weight & Relative choice wt & Normal & +0.48 & [0.38, 0.58] & Mandeville et al. (2014) & C \\
\texttt{PAR\_BETA\_INFRA\_DCE} & Infra DCE Weight & Relative choice wt & Normal & +0.36 & [0.27, 0.45] & Mandeville et al. (2014) & C \\
\texttt{PAR\_BETA\_RURAL\_DCE} & Rural DCE Weight & Relative choice wt & Normal & +0.31 & [0.22, 0.40] & Mandeville et al. (2014) & C \\
\texttt{PAR\_SCEN1\_WAGE\_INC} & Salary Increment & Wage policy & Uniform & 0.30 & [0.20, 0.50] & Scenario 1 Specification & E \\
\texttt{PAR\_SCEN1\_RRR\_WAGE} & Wage Retention & Policy RRR & LogNorm & 0.12 & [0.08, 0.18] & Wage Elasticity Meta-analysis & C \\
\texttt{PAR\_SCEN2\_RRR\_INFRA} & Infra Retention & Policy RRR & LogNorm & 0.10 & [0.06, 0.16] & Scenario 2 Specification & E \\
\texttt{PAR\_SCEN3\_RRR\_RURAL} & Rural Retention & Policy RRR & LogNorm & 0.18 & [0.12, 0.26] & Scenario 3 Specification & E \\
\texttt{PAR\_SCEN4\_QUOTA\_EXP} & Quota Increase & Inflow multiplier & Uniform & 0.15 & [0.10, 0.25] & Scenario 4 Specification & E \\
\texttt{PAR\_SCEN4\_ROS\_COMP} & Return-of-Service & Service adherence & Beta & 0.75 & [0.65, 0.85] & B{\"a}rnighausen \& Bloom (2009) & C \\
\texttt{PAR\_POLICY\_DECAY} & Policy Decay & Efficacy loss & Beta & 0.035 & [0.02, 0.06] & Health Economics Standard & D \\
\texttt{PAR\_COMP\_COST\_FAC} & Clinical Support & Support multiplier & Gamma & 0.35 & [0.25, 0.45] & CHEERS 2022 Item 13 & C \\
\texttt{PAR\_DISCOUNT\_RATE} & Social Discount & Annual discount & Discrete & 0.03 & [0.00, 0.08] & Second Panel Standard & D \\
\texttt{PAR\_GROSS\_TRAIN\_MD} & Gross Training Cost & Training outlay & Gamma & 85,000 & [65k, 110k] & MESRS Educational Accounts & A \\
\texttt{PAR\_RES\_LABOR\_OFF} & Resident Labor Offset & Clinical labor & Beta & 0.40 & [0.30, 0.50] & Human Capital Accounting & D \\
\texttt{PAR\_NET\_TRAIN\_COST} & Net Training Loss & Net fiscal cost & Gamma & 51,000 & [39k, 66k] & Net Fiscal Training Loss & D \\
\texttt{PAR\_BASE\_SALARY\_ANN}& Base Public Salary & Annual wage & Gamma & 36,000 & [30k, 48k] & Statut Fonction Publique & A \\
\texttt{PAR\_INFRA\_CAPEX\_ANN}& Hospital Capex & Annual capital & Gamma & 25.0M & [15M, 40M] & MOH Capital Envelope & E \\
\bottomrule
\end{tabular}%
}
\end{table}

\newpage

\subsection{Appendix C: CHEERS 2022 Reporting Checklist}

\begin{table}[H]
\centering
\scriptsize
\caption{Consolidated Health Economic Evaluation Reporting Standards (CHEERS 2022) Checklist.}
\label{tbl:app_cheers}
\resizebox{\textwidth}{!}{%
\begin{tabular}{cp{5.5cm}cp{8cm}}
\toprule
\textbf{Item} & \textbf{CHEERS Item Name} & \textbf{Reported} & \textbf{Manuscript Location \& Description} \\
\midrule
1 & Title & Yes & Identifies dynamic workforce simulation, policy counterfactuals, and economic evaluation. \\
2 & Abstract & Yes & Structured summary with background, methods, results, and policy conclusions. \\
3 & Background and Objectives & Yes & Section 1: Institutional context, emigration surge, and explicit study aims. \\
4 & Health Economic Analysis Plan & Yes & Section 2: CHEERS 2022 and Second Panel compliant plan specified prior to execution. \\
5 & Study Population & Yes & Section 2: Entire licensed and practicing physician workforce in Tunisia (36-cell tensor). \\
6 & Setting and Location & Yes & Section 2: National Tunisian health system (Coastal vs. Interior macro-regions). \\
7 & Comparators & Yes & Section 2: Status Quo baseline (SCEN\_0) versus 5 standardized policy counterfactuals. \\
8 & Perspective & Yes & Section 2: Public payer perspective (Ministry of Health and Ministry of Finance). \\
9 & Time Horizon & Yes & Section 2: 25-year projection horizon (2025--2050) with milestones at 2030, 2040, and 2050. \\
10 & Discount Rate & Yes & Section 2: Baseline 3.0\% annual discount rate; tested at 0.0\%, 5.0\%, and 8.0\%. \\
11 & Selection of Outcomes & Yes & Section 2: Public and interior physician-years retained; emigrants averted. \\
12 & Measurement of Outcomes & Yes & Section 2: Continuous-time competing-risk multi-state Markov transitions. \\
13 & Valuation of Outcomes & Yes & Section 2: Natural units (physician-years); complementary operational support factor (35\%). \\
14 & Resource Use and Costs & Yes & Section 2: Inframarginal wage bill expansion, capital budgets, and net training losses. \\
15 & Currency and Price Date & Yes & Section 2: 2024 Tunisian Dinars (TND); deflated by INS Consumer Price Index. \\
16 & Model Type & Yes & Section 2: Discrete-time 36-cell Markov model with continuous-time hazard engine. \\
17 & Assumptions & Yes & Section 2 \& Appendix A: Documented baseline filter ($\alpha$), realization ($\kappa$), and decay. \\
18 & Analytical Methods & Yes & Section 2: 2-parameter MLE calibration, temporal validation split, and external validation. \\
19 & Baseline Results & Yes & Section 3: Calibrated 36-cell stock, in-sample and out-of-sample fit metrics (MAPE, RMSE). \\
20 & Resource Use Results & Yes & Section 3 \& Table 3: Undiscounted and discounted fiscal costs across all scenarios. \\
21 & Cost-Effectiveness Results & Yes & Section 3: CCA, ICERs ($\text{ICER}_{\text{public}}, \text{ICER}_{\text{interior}}$), and efficiency frontier. \\
22 & Uncertainty Analysis & Yes & Section 3: Deterministic Tornado analysis and 1,000-draw Monte Carlo PSA. \\
23 & Heterogeneity & Yes & Section 3: Disaggregated outcomes by career stage, sector, geography, and specialty. \\
24 & Discussion: Findings & Yes & Section 4: Training-only trap, infrastructure efficiency, and fiscal feasibility. \\
25 & Discussion: Limitations & Yes & Section 4: Aggregated specialty clusters, stationary return migration, and lack of DALYs. \\
26 & Generalizability & Yes & Section 4: Transferability of findings to Maghreb and middle-income health systems. \\
27 & Funding & Yes & Declarations: Independent academic research without external commercial sponsorship. \\
28 & Conflicts of Interest & Yes & Declarations: Authors declare no competing financial or professional interests. \\
\bottomrule
\end{tabular}%
}
\end{table}

\newpage

\subsection{Appendix D: STROBE Statement Checklist}

\begin{table}[H]
\centering
\scriptsize
\caption{STROBE Statement Checklist for Observational and Administrative Data Components.}
\label{tbl:app_strobe}
\resizebox{\textwidth}{!}{%
\begin{tabular}{cp{5.5cm}cp{8cm}}
\toprule
\textbf{Item} & \textbf{STROBE Checklist Item} & \textbf{Reported} & \textbf{Manuscript Location \& Description} \\
\midrule
1 & Title and Abstract & Yes & Identifies administrative census and migration registry time series. \\
2 & Background / Rationale & Yes & Section 1: Explains scientific rationale for examining Tunisian physician emigration. \\
3 & Objectives & Yes & Section 1: Specifies workforce projection and empirical calibration aims. \\
4 & Study Design & Yes & Section 2: Multi-source administrative time-series (2015--2024) synthesis. \\
5 & Setting & Yes & Section 2: National Tunisian health system and primary destination countries. \\
6 & Participants / Eligibility & Yes & Section 2: All licensed physicians practicing in Tunisia or requesting clearance. \\
7 & Variables & Yes & Section 2: Headcount stock, sector, governorate, specialty, and departure flows. \\
8 & Data Sources & Yes & Section 2: INS Stat 123, Carte Sanitaire 2023, CNOM, French DREES, German B{\"A}K. \\
9 & Bias & Yes & Section 2: Active practice gap corrected via $\alpha = 0.88$; clearance via $\lambda_{\text{exit}} = 0.75$. \\
10 & Study Size & Yes & Section 2: Complete national census of practicing physicians ($N = 15,315$). \\
11 & Quantitative Variables & Yes & Section 2: Continuous headcounts, densities per 10k, and costs in constant TND. \\
12 & Statistical Methods & Yes & Section 2: IPF raking, Maximum Likelihood calibration, and temporal holdout testing. \\
13 & Descriptive Data & Yes & Section 3: Baseline distribution across 36 demographic cells presented in Table 1. \\
14 & Outcome Data & Yes & Section 3: Annual realized emigration flows and projected active stocks (2025--2050). \\
15 & Main Results & Yes & Section 3: Observed historical fits (MAPE = 1.61\% stock) and policy milestones. \\
16 & Other Analyses & Yes & Section 3: Parameter tornado swings and 1,000-draw Monte Carlo PSA. \\
17 & Key Results & Yes & Section 4: Synthesis of historical trends and public-sector depletion. \\
18 & Limitations & Yes & Section 4: Data gaps, returnee registry absence, and intention-to-action gaps. \\
19 & Interpretation & Yes & Section 4: Balanced interpretation of supply expansion vs. retention. \\
20 & Generalizability & Yes & Section 4: Applicability of findings to Maghreb and North African health systems. \\
21 & Funding & Yes & Declarations: Non-commercial academic research declaration. \\
\bottomrule
\end{tabular}%
}
\end{table}

\newpage

\subsection{Appendix E: PRISMA Evidence Synthesis Workflow}

\begin{table}[H]
\centering
\scriptsize
\caption{PRISMA Evidence Synthesis Matrix and Search Governance.}
\label{tbl:app_prisma}
\resizebox{\textwidth}{!}{%
\begin{tabular}{p{3.5cm}p{6.5cm}p{4cm}p{4cm}}
\toprule
\textbf{PRISMA Domain} & \textbf{Screening Criteria \& Protocol} & \textbf{Included Evidence in Registry} & \textbf{Parameter Role} \\
\midrule
\textbf{Tunisian Surveys} & Peer-reviewed cross-sectional studies reporting emigration intention and odds ratios. & Ben Romdhane et al. (2022), Dhouib et al. (2024), Ben Salem et al. (2024). & Baseline intention ($p_0 = 0.686, 0.756$), survey ORs converted to RRs ($\text{RR} = 1.14$--$1.23$). \\
\textbf{Maghreb Evidence} & Comparative physician migration surveys from Morocco, Algeria, Egypt, and Lebanon. & Sylla et al. (2021), El Morabety (2022), Gouda et al. (2021), Akl et al. (2008). & Regional validation benchmarks for intention prevalence and diaspora shares (Class B). \\
\textbf{DCE Meta-Analyses} & Systematic reviews and pooled meta-analyses of discrete choice experiments in LMICs. & Mandeville et al. (2014), Lagarde \& Blaauw (2009). & Relative choice utility weights for wages ($\beta_{\text{wage}} = 0.48$), infra ($\beta_{\text{infra}} = 0.36$), and rural ($\beta_{\text{rural}} = 0.31$). \\
\textbf{Realization Panels} & Multi-country microdata panels evaluating stated intentions vs. realized migration. & Clemens \& Mendola (2020), Docquier et al. (2014), Creighton (2013). & Prior distributions for Realization Factor ($\kappa \sim \text{Beta}(35, 65)$, base $\kappa = 0.35$). \\
\textbf{Administrative Data} & Official national census reports, ministerial yearbooks, and foreign licensing registers. & INS Stat 123, MOH Carte Sanitaire 2023, CNOM, French DREES, German B{\"A}K. & Baseline 36-cell stock tensor, outflow calibration targets, active ratio ($\alpha = 0.88$). \\
\bottomrule
\end{tabular}%
}
\end{table}

\newpage

\bibliographystyle{plainnat}
\bibliography{references}

\end{document}
